\section{Integration Part 1}
\tsDef{Antiderivative / Indefinite Integral}{
Let $I\subseteq\mathbb{R}$ be an interval and let
\(
f:I\to\mathbb{R}
\)
be a function.
A differentiable function
\(
F:I\to\mathbb{R}
\)
satisfying
\(
F'(x)=f(x)
\quad \forall x\in I
\)
is called an \emph{antiderivative} of $f$ on $I$.
Determining $F$ from $f$ is called \emph{indefinite integration}.
}
\\
\tsIdea{Non-Uniqueness of Antiderivatives}{
An antiderivative is not unique.
If $F(x)$ is an antiderivative of $f(x)$, then
\(
F(x)+C
\)
is also an antiderivative for every constant $C\in\mathbb{R}$.
Hence a function possesses an entire family of antiderivatives.
This family is called the \emph{indefinite integral} and is written as
\[
\int f(x)\,dx.
\]
The function $f$ is called the \emph{integrand}.
}
\\
\tsThe{Linearity of the Indefinite Integral}{
For indefinite integrals:
\[
\int (f(x)\pm g(x))\,dx
=
\int f(x)\,dx
\pm
\int g(x)\,dx
\]
and
\[
\int s f(x)\,dx
=
s\int f(x)\,dx,
\]
where $s$ is a constant.
}
\\
\tsThe{Basic Integration Rules}{
Rules:
\begin{itemize}[noitemsep]
    \item \(
    \int k\,dx = kx + C
    \)

    \item \(
    \int x^n\,dx
    =
    \frac{1}{n+1}x^{n+1}+C,
    \qquad n\neq -1
    \)

    \item \(
    \int e^x\,dx = e^x + C
    \)

    \item \(
    \int \frac{1}{x}\,dx
    =
    \ln(|x|)+C
    \)

    \item \(
    \int \sin(x)\,dx
    =
    -\cos(x)+C
    \)

    \item \(
    \int \cos(x)\,dx
    =
    \sin(x)+C
    \)

    \item \(
    \int \sinh(x)\,dx
    =
    \cosh(x)+C
    \)

    \item \(
    \int \cosh(x)\,dx
    =
    \sinh(x)+C
    \)
\end{itemize}
}
\tsThe{Integration by Parts}{
Using the product rule:
\(
\frac{d}{dx}(f(x)g(x))
=
f'(x)g(x)+f(x)g'(x),
\)
we obtain
\(
\int \frac{d}{dx}(f(x)g(x))\,dx
=
(fg)(x).
\)
Therefore,
\(
(fg)(x)
=
\int f'(x)g(x)\,dx
+
\int f(x)g'(x)\,dx.
\)
Rearranging gives the formula for \emph{integration by parts}:
\[
\boxed{
\int f'(x)g(x)\,dx
=
f(x)g(x)
-
\int f(x)g'(x)\,dx
}
\]
}
\\
\tsThe{Integration by Substitution}{
Using the chain rule:
\(
\frac{d}{dx}f(g(x))
=
f'(g(x))g'(x),
\)
we obtain
\(
\int f'(g(x))g'(x)\,dx
=
f(g(x))+C.
\)
With the substitution
\(
y=g(x),
\)
this becomes
\(
\int f'(g(x))g'(x)\,dx
=
\int \frac{df}{dy}\,dy.
\)
Hence the substitution rule is
\[
\boxed{
\int f'(g(x))g'(x)\,dx
=
f(g(x))+C
}
\]
}
\section{Integration Part 2}
% \tsIdea{Motivation}{
% Given a velocity function $v(t)$ describing one-dimensional motion, natural questions are:
% \begin{itemize}
% \item What is the total distance traveled?
% \item What is the change in position?
% \end{itemize}
% }
\tsDef{Partition of an Interval}{
Let $[a,b]\subseteq\mathbb{R}$ be a compact interval.
A \emph{partition} of $[a,b]$ is a finite set of points
\[
a=x_0<x_1<\cdots<x_{n-1}<x_n=b.
\]
The points $x_i$ are called \emph{subdivision points}.
A partition
\[
a=y_0<y_1<\cdots<y_m=b
\]
is called a \emph{refinement} of
\(
a=x_0<x_1<\cdots<x_n=b
\)
if
\(
\{x_0,\dots,x_n\}\subseteq\{y_0,\dots,y_m\}.
\)
}
\\
\tsDef{Step Function}{
A function
\(
f:[a,b]\to\mathbb{R}
\)
is called a \emph{step function} if there exists a partition
\(
a=x_0<x_1<\cdots<x_n=b
\)
such that for every $k=1,\dots,n$, the restriction
\(
f|_{(x_{k-1},x_k)}
\)
is constant.
}
\\
\tsDef{Integral of a Step Function}{
Let
\(
f:[a,b]\to\mathbb{R}
\)
be a step function with respect to the partition
\(
a=x_0<x_1<\cdots<x_n=b.
\)
Then the integral of $f$ over $[a,b]$ is defined by
\[
\int_a^b f(x)\,dx
:=
\sum_{k=1}^n
c_k(x_k-x_{k-1}),
\]
where $c_k$ denotes the constant value of $f$ on $(x_{k-1},x_k)$.
}
\\
\tsDef{Upper and Lower Sums}{
Let
\[
f:[a,b]\to\mathbb{R}.
\]
The set of upper sums is defined by
\[
U(f)
:=
\left\{
\int_a^b s(x)\,dx
\;\middle|\;
s\in TF,\ s\ge f
\right\},
\]
and the set of lower sums by
\[
L(f)
:=
\left\{
\int_a^b s(x)\,dx
\;\middle|\;
s\in TF,\ s\le f
\right\},
\]
where $TF$ denotes the set of step functions.
}
\\
\tsDef{Riemann Integral}{
A bounded function
\(
f:[a,b]\to\mathbb{R}
\)
is called \emph{Riemann integrable} if
\[
\sup L(f)=\inf U(f).
\]
In this case, the common value is called the \emph{Riemann integral} and is denoted by
\[
\int_a^b f(x)\,dx
:=
\sup L(f)=\inf U(f).
\]
The numbers $a,b$ are called the lower and upper limits of integration, and $f$ is called the integrand.
}
\\
\tsThe{Properties of the Riemann Integral}{
    Properties of the Riemann Integral:
    \begin{itemize}
        \item[(i)] (Linearity)
        If $f,g:[a,b]\to\mathbb{R}$ are integrable and $\alpha,\beta\in\mathbb{R}$, then
        \[
        \int_a^b (\alpha f+\beta g)(x)\,dx
        =
        \alpha\int_a^b f(x)\,dx
        +
        \beta\int_a^b g(x)\,dx.
        \]

        \item[(ii)] (Monotonicity)
        If
        \(
        f(x)\le g(x)
        \)
        for all $x$, then
        \[
        \int_a^b f(x)\,dx
        \le
        \int_a^b g(x)\,dx.
        \]

        \item[(iii)] (Triangle Inequality)
        \[
        \left|
        \int_a^b f(x)\,dx
        \right|
        \le
        \int_a^b |f(x)|\,dx.
        \]

        \item[(iv)] (Reversal of Integration Limits)
        \[
        \int_a^b f(x)\,dx
        =
        -\int_b^a f(x)\,dx.
        \]

        \item[(v)] (Splitting the Interval)
        For $a\le c\le b$:
        \[
        \int_a^b f(x)\,dx
        =
        \int_a^c f(x)\,dx
        +
        \int_c^b f(x)\,dx.
        \]
    \end{itemize}
}

\tsThe{Monotone Functions are Integrable}{
Every monotone function
\(
f:[a,b]\to\mathbb{R}
\)
is Riemann integrable.
}
\\
\tsThe{Continuous Functions are Integrable}{
Every continuous function
\(
f:[a,b]\to\mathbb{R}
\)
is Riemann integrable.
}
\\
\tsDef{Pointwise Convergence of Functions}{
Let
\(
(f_n)_{n\in\mathbb{N}_0}
\)
be a sequence of functions
\(
f_n:D\to\mathbb{R},
\)
and let
\(
f:D\to\mathbb{R}.
\)
The sequence $(f_n)$ converges \emph{pointwise} to $f$ if for every $x\in D$, the sequence
\(
(f_n(x))_{n\in\mathbb{N}_0}
\)
converges to $f(x)$.
}
\\
\tsDef{Uniform Convergence}{
The sequence $(f_n)$ converges \emph{uniformly} to $f$ if
\(
\forall \varepsilon>0\ \exists N
\)
such that for all
\(
n\ge N
\ \text{and all }x\in D,
\)
we have
\(
|f_n(x)-f(x)|<\varepsilon.
\)
}
\\
\tsThe{Uniform Limit of Continuous Functions}{
Let
\(
(f_n)_{n\in\mathbb{N}_0}
\)
be a sequence of continuous functions
\(
f_n:D\to\mathbb{R}
\)
converging uniformly to
\(
f:D\to\mathbb{R}.
\)
Then $f$ is continuous.
}
\\
\tsThe{Uniform Limit and Integration}{
Let
\(
(f_n)_{n\in\mathbb{N}_0}
\)
be a sequence of integrable functions
\(
f_n:[a,b]\to\mathbb{R}
\)
converging uniformly to
\(
f:[a,b]\to\mathbb{R}.
\)
Then $f$ is integrable and
\[
\int_a^b f(x)\,dx
=
\lim_{n\to\infty}
\int_a^b f_n(x)\,dx.
\]
}
\\
\tsThe{Mean Value Theorem for Integrals}{
If
\(
f:[a,b]\to\mathbb{R}
\)
is continuous, then there exists
\(
c\in[a,b]
\)
such that
\(
f(c)
=
\frac{1}{b-a}
\int_a^b f(x)\,dx.
\)
The quantity
\[
\frac{1}{b-a}
\int_a^b f(x)\,dx
\]
is called the \emph{average value} of $f$ on $[a,b]$.
}
\\
\tsThe{Fundamental Theorem of Calculus}{
Let
\(
f:[a,b]\to\mathbb{R}
\)
be continuous.
Then for every constant $C\in\mathbb{R}$, the function
\[
F(x)
:=
\int_a^x f(y)\,dy + C
\]
is an antiderivative of $f$.
Moreover, every antiderivative of $f$ is of this form.
}
\\
\tsCor{Integral Representation}{
Let
\(
F:[a,b]\to\mathbb{R}
\)
be continuously differentiable.
Then
\[
F(x)
=
F(a)
+
\int_a^x F'(t)\,dt.
\]
}
\\
\tsCor{Newton--Leibniz Formula}{
Let
\(
f:[a,b]\to\mathbb{R}
\)
be continuous and let $F$ be an antiderivative of $f$.
Then
\[
\int_a^b f(t)\,dt
=
F(b)-F(a).
\]
}
\\
\tsThe{Integration by Parts}{
Let
\(
f,g:[a,b]\to\mathbb{R}
\)
be continuously differentiable.
Then
\[
\int_a^b f(x)g'(x)\,dx
=
f(x)g(x)\Big|_a^b
-
\int_a^b f'(x)g(x)\,dx.
\]
}
\\
\tsThe{Substitution Rule (Version 1)}{
Let
\(
f:I\to J
\)
be continuously differentiable and
\(
g:J\to\mathbb{R}
\)
continuous.
Then for every
\(
[a,b]\subseteq I,
\)
\[
\int_a^b g(f(x))f'(x)\,dx
=
\int_{f(a)}^{f(b)} g(y)\,dy.
\]
}
\\
\tsThe{Substitution Rule (Version 2)}{
Let
\(
f:I\to J
\)
be continuously differentiable and
\(
g:J\to\mathbb{R}
\)
continuous.
Assume
\(
f'(x)\neq 0
\quad\forall x\in[a,b]
\)
and let
\(
f^{-1}:[f(a),f(b)]\to\mathbb{R}
\)
denote the inverse of $f|_{[a,b]}$.
Then
\[
\int_a^b g(f(x))\,dx
=
\int_{f(a)}^{f(b)}
g(y)(f^{-1})'(y)\,dy.
\]
}